\specialchapter{第 IV、V、VI 章导言}

对半单（解析的或代数的）群及其李代数的研究，引出了对根系、Coxeter 群和 Tits 系统的考察。第 IV、V、VI 章专门讨论这些结构。

为了帮助读者了解方向，下面给出几个例子。

I. 设 g 为一个复半单李代数，h 为 \(g^{1}\) 的一个 Cartan 子代数。关于 h 的 g 的一个根，是 h 上的一个非零线性形式 \(\alpha\)，满足存在 g 的一个非零元素 x，使得 \([h, x] = \alpha(h)x\) 对所有 \(h \in h\) 成立。这些根在与 h 对偶的向量空间 \(h^{*}\) 中构成一个约化根系 R。给定 R，便在同构意义下确定 g，并且每个约化根系都同构于以此方式得到的一个根系。g 的保持 h 稳定的一个自同构确定 \(h^{*}\) 的一个保持 R 不变的自同构，而 R 的每个自同构都以此方式得到。R 的 Weyl 群由 \(h^{*}\) 的所有自同构组成，这些自同构由 g 的保持 h 稳定的内自同构所确定；该群是一个 Coxeter 群。

设 \(G\) 为一个连通复李群，其李代数为 \(\mathfrak{g}\)，并设 \(\Gamma\) 为 \(G\) 中的元素 \(h\) 组成的子群，这些元素满足 \(\exp_G(2\pi ih) = 1\)。设 \(R^-\) 为 \(\mathfrak{h}\) 中与 \(R\) 互为逆的根系，设 \(Q(R^-)\) 为 \(\mathfrak{h}\) 中由 \(R^-\) 生成的子群，并设 \(P(R^-)\) 为与子群 \(Q(R)\)（\(\mathfrak{h}^*\) 的子群，由 \(R\) 生成）相关联的子群（即由 \(h \in \mathfrak{h}\) 组成，满足 \(\lambda(h)\) 对所有 \(\lambda \in Q(R)\) 都是整数）。于是 \(P(R^-) \supset \Gamma \supset Q(R^-)\)。此外，\(G\) 的中心与 \(P(R^-)/\Gamma\) 有典范同构，而 \(G\) 的基本群与 \(\Gamma/Q(R^-)\) 有典范同构。特别地，\(\Gamma\) 等于 \(P(R^-)\) 当 \(G\) 是伴随群时，并且 \(\Gamma\) 等于 \(Q(R^-)\) 当 \(G\) 是单连通时。最后，有限维线性表示（关于 \(G\)）的权是 \(\mathfrak{h}^*\) 的子群中与 \(\Gamma\) 相关联的元素。

\begin{enumerate}
\def\labelenumi{\Roman{enumi}.}
\setcounter{enumi}{1}
\tightlist
\item
  设 G 为一个半单连通紧实实李群，g 为其李代数。设 T 为 G 的极大环面，其李代数为 t，并设 X 为 T 的特征标群。设 R 为 X 中满足下述条件的非零元素 \(\alpha\) 的集合：存在 g 的非零元素 x，使得 (Ad t). \(x = \alpha(t)x\) 对所有 \(t \in T\) 成立。将 X 视为实向量空间 \(V = X \otimes_{Z} R\) 中的一个格；于是 R 是 V 中的约化根系。设 N 为 G 中 T 的规范化子；N 在 T 上的作用给出从群 N/T 到 R 的 Weyl 群的一个同构。我们有 \(P(R) \supset X \supset Q(R)\)；此外，当 G 是单连通时 \(X = P(R)\)，当 G 的中心退化为单位元时 \(X = Q(R)\)。
\end{enumerate}

g 的复化李代数 \(\mathfrak{g}_{(\mathbf{C})}\) 是半单的，而 \(t_{(\mathbf{C})}\) 是它的一个 Cartan 子代数。存在从 \(V_{(\mathbf{C})}\) 到 \(t_{(\mathbf{C})}\) 的对偶空间的一个典范同构，它将 R 变为 \(\mathfrak{g}_{(\mathbf{C})}\) 关于 \(t_{(\mathbf{C})}\) 的根系。

\begin{enumerate}
\def\labelenumi{\Roman{enumi}.}
\setcounter{enumi}{2}
\tightlist
\item
  设 G 为交换域 k 上的连通半单代数群。设 T 是 G 的在 k 上分裂的环面集合中的极大元，并设 X 为 T 的特征标群（即从 T 到乘法群的同态）。将 X 视为实向量空间 \(V = X \otimes_{Z} R\) 中的一个格。G 关于 T 的根，是 X 中满足下述条件的非零元素 \(\alpha\)：存在 G 的李代数 g 的非零元素 x，使得 \((\operatorname{Ad} t)x = \alpha(t)x\) 对所有 \(t \in T\) 成立。这在 V 中给出一个根系 R，但它不一定是约化的。设 N 为 G 中 T 的规范化子，Z 为 T 在 G 中的中心化子，并设 \(N(k)\) 和 \(Z(k)\) 为它们在 k 上的有理点群。\(N(k)\) 在 T 上的作用给出从 \(N(k)/Z(k)\) 到 R 的 Weyl 群的一个同构。
\end{enumerate}

设 U 为 G 的幂零子群集合中的极大元，它在 k 上定义并由 Z 规范化。置 P = Z.U。则 \(P(k) = Z(k)\) 。\(U(k)\) 且 \(P(k) \cap N(k) = Z(k)\) 。此外，存在 R 的一个基 \((\alpha_{1}, \ldots, \alpha_{n})\)，使得 T 在 U 中的权是关于此基的 R 的正根；若 S 表示 \(N(k)/Z(k)\) 中通过上述同构对应于对称 \(s_{\alpha_{i}} \in W(R)\)（这些对称与根 \(\alpha_{i}\) 相关）的元素集合，则四元组 \((G(k), P(k), N(k), S)\) 是一个 Tits 系统。

\begin{enumerate}
\def\labelenumi{\Roman{enumi}.}
\setcounter{enumi}{3}
\tightlist
\item
  在局部域上的半单代数群理论中，会遇到这样的 Tits 系统，其 Weyl 群是某个根系的仿射 Weyl 群。例如，设 \(\mathbf{G} = \mathbf{S}\mathbf{L}(n + 1, \mathbf{Q}_{p})\)（其中 \(n \geq 1\)）。设 B 为由矩阵 \((a_{ij}) \in \mathbf{S}\mathbf{L}(n + 1, \mathbf{Z}_{p})\) 组成的群，这些矩阵满足 \(a_{ij} \in pZ_{p}\)，其中 i \textless{} j，并设 N 为 G 的子群，它由每行每列恰有一个非零元素的矩阵组成。于是存在 \(N/(B \cap N)\) 的一个子集 S，使四元组 \((\mathbf{G}, \mathbf{B}, \mathbf{N}, \mathbf{S})\) 成为一个 Tits 系统。群 \(W = N/(B \cap N)\) 是一个 \(A_{n}\) 型根系的仿射 Weyl 群；这是一个无限 Coxeter 群。
\end{enumerate}

在准备这些章节的过程中，我们多次与 J. Tits 交谈，获益良多。谨向他致以衷心的感谢。
